% Source: 29_Chapter_20_The_Nectar_of_Learning_Scholasticism.txt
\bookchapter[The Nectar of Learning]{CHAPTER 20}{THE NECTAR OF LEARNING}{Scholasticism}

\leadin{F}{OR EIGHT CENTURIES GOTHIC} cathedrals throughout Europe have inspired worshipers, and
awed tourists. The medieval masters of Gothic style tried to portray in stone and glass man’s
central religious quest. They wanted to depict a tension: on the one hand was man aspiring to reach
the heights of heaven; on the other hand was God condescending to address the least of men.

The movement of the Gothic, therefore, is two-way. The pillars, arches, and steeples, aligned like
rows of rockets ready to ascend to heaven, point skyward. But through colorful windows of leaded
glass the light of God descends to meet the lowly. It is an architect’s version of human reason and
divine revelation.

Windows in the Gothic structure could be plentiful and large because the outside pillars and
buttresses eliminated the need for thick walls. Between the shafts of slender stone, architects
could set windows of colored glass. Artists used brilliant colors—crimson, blue, purple, ruby—to
tell the story of redemption from creation to the Last Judgment. Worshipers, then, could ponder the
sacrifice of Isaac, the crossing of the Red Sea, the birth of Jesus, or Saint Anthony’s struggle
with devils.

The total effect was breathtaking. As the evening sun cast its warm enchanting rays upon the cold
gray stone, even the cobbler from his pew could sense his kinship with Moses, Isaiah, Jesus, Paul,
Augustine, and Benedict.

Light, which played upon the multicolored windows, had been a constant symbol of God and his ways
with men. “God is light,” wrote the apostle. “In him is no darkness at all” (1 John 1:5). And Jesus
himself used this image of his own incarnation when he announced, “I am the light of the world.”

The Gothic cathedral, therefore, displayed the spiritual tension of the Christian drama: the highest
aspirations of man and the descending light of God. Man, in effect, ascends while God descends. Such
language is, of course, figurative. God is no more above than below in any spatial sense. But man
has always described his need in terms of reach, and God’s truth in terms of descent.

\section{THE RISE OF UNIVERSITIES}

Appropriately, therefore, schools in these cathedrals gave birth to medieval universities, for the
supreme task of the university was to understand and explain the light of God’s revealed truth. Just
as the crusades reflected a widespread passion to extend the authority of God in Muslim lands, so
the universities reveal an intense hunger to understand the truth of God received from any land.

But just how did the world of ideas bow to the rule of God? How was reason made the servant of
faith? We call this period in the history of Christian thought “Scholasticism” because a distinctive
method of scholarship arose and because a unique theology of the Middle Ages emerged. The aim of the
Schoolmen, as these teachers are sometimes called, was twofold: to reconcile Christian doctrine and
human reason, and to arrange the teachings of the church in an orderly system.

A free search for the truth was never in view since the chief doctrines of the Christian faith were
regarded as fixed. The purpose of discussion was to show the reasonableness of the doctrines and to
explain their implications. The whole endeavor surged forward on the winds of intellectual
curiosity. One eleventh-century scholar from Liege represents hosts of others:

\begin{blockquote}
Olbert was not able to satiate his thirst for study. When he would hear of someone distinguished in
the arts he flew there at once, and the more he thirsted the more he absorbed something delightful
from each master~.~.~.~. Afterwards just like the bees among flowers, gorged with the nectar of
learning, he returned to the hive and lived there studiously in a religious way, and religiously in
a studious manner.
\end{blockquote}

Living “religiously in a studious manner” was the hallmark of medieval education. Its central aim
was eternal salvation. Originally, however, learning had been limited primarily to the clergy.
Benedict of Nursia had insisted that his monks read and study for their own spiritual development.
As a result, Benedictine abbeys created a kind of rudimentary schooling in Latin.

Later, in the eighth century, Charlemagne, dreaming of a Christian empire, had widened the
opportunities for study through a decree that every monastery have a school to teach all those “who
with God’s help are able to learn.” The emperor himself set an example with a palace school for his
own children and those of his courtiers.

The best chance for learning among laymen came from cathedral schools. Since cathedrals, churches of
the bishops, were located in towns, their schools to train parish priests were in time opened to
all.

The curriculum of the cathedral school was limited to grammar, rhetoric, logic, arithmetic,
geometry, music, and astronomy—the seven liberal arts, so called because in ancient Rome their study
had been reserved for liberi, freemen. The few texts available were writings of a handful of
scholars of the early Middle Ages. One was Cassiodorus, a sixth-century Roman whose Handbook of
Sacred and Secular Learning defined the liberal arts and interpreted the Bible. Another was
Boethius, a contemporary of Cassiodorus, whose Consolation of Philosophy, written while he was in
prison for an alleged plot against the barbarian King Theodoric, attempted to reconcile the
misfortunes of man with the concept of a benevolent, omnipotent God. These men, along with
Augustine, Pope Gregory the Great, and a few other church fathers, were known as auctores,
authorities, whose words the medieval student would not lightly presume to contradict.

\section{THE MAGNETISM OF AN ABLE TEACHER}

A new day dawned, however, with the coming of the great schoolmasters. We can trace the birth of
universities to the magnetism of popular teachers whose skill and enthusiasm for learning attracted
students wherever they happened to be.

As Anne Fremantle describes them, the first of this new breed was Gerbert, master of the cathedral
school at Rheims in the latter half of the tenth century. The son of a serf, Gerbert was destined to
end his days as Pope Sylvester II (999–1003). But he left a deeper mark upon history through his
scholarship. “As a young monk, Gerbert had been so brilliant a student that his abbot had taken the
unusual step of sending him to Spain to study mathematics. Although Gerbert’s mentor there was a
Christian bishop, he was also exposed to the broad and tolerant culture of the ruling Muslims.” This
was the first of a number of significant contributions Muslims made to the Christian intellectual
awakening.

Gerbert returned to Rheims deeply impressed by the inquisitive, questing spirit of Muslim
scholarship. When he began to teach at Rheims, he announced that quotations of traditional church
authorities no longer sufficed; his pupils were henceforth to study Roman classics in the original.
“To this end Gerbert collected manuscripts wherever he could and built a sizable library—no mean
feat at a time when a manuscript sometimes took a year to copy, and cost at least the equivalent of
a churchman’s yearly income.”

The liveliest figure in this early stage of intellectual revolution was Peter Abelard (1079–1142).
The eldest son of a minor noble of Brittany (northwestern France), “Abelard for love of learning had
given up his inheritance rights to younger brothers, and roamed France to sit at the feet of the
great masters, now listening, now openly challenging them in class.” In time he established himself
as a lecturer in Paris, where he attracted a host of students. He also began to write.

In a treatise entitled Sic et Non (Yes and No), he posed 158 questions from Christian teaching and
answered them with conflicting quotations from the Scriptures, the church fathers, and pagan
classics. “The first key to wisdom,” Abelard asserted, “is assiduous and frequent
questioning~.~.~.~. For by doubting we come to inquiry, and by inquiry we arrive at the truth.” This
idea, commonplace to the Greeks, was hardly so to medieval Europeans.

Abelard’s zest for doubt won the applause of some, but alarmed as many others. Another of his books,
on the nature of the Trinity, invoked condemnation by a church council at Soissons in 1121, and the
brilliant scholar found himself behind the walls of a monastery.

Resourceful as usual, Abelard secured the monastery’s permission to leave a year later to live in
the wilderness southeast of Paris. “Students flocked to his side. They built him a shelter, tilled
his land, and begged him to teach once more. Resuming his pursuit of reason, Abelard again and again
fell afoul of conservatives in the church, this time including Abbot Bernard of Clairvaux, the most
influential churchman in Christendom.” Bernard pursued Abelard as devoutly as he preached the Second
Crusade. “The faith of the righteous believes,” he declared, “it does not dispute.” At Bernard’s
instigation, a church council at Sens in 1140 condemned Abelard for heresy. Abelard retired to the
abbey of Cluny, where he stayed in seclusion for the remaining two years of his life.

No one, however, could stifle the growth of the seeds he had scattered. Schools sprouted all around
the Continent. Fewer than one hundred years after his death, universities flourished at Paris,
Orleans, and Montpellier in France; across the English Channel at Oxford and Cambridge; at Bologna
and Padua in Italy.

The event that marked the flowering of the universities was the grouping of students and masters
into guilds. “As craftsmen had done before them,” explains Fremantle, scholars banded together for
mutual interest and protection, and called themselves a universitas, the medieval name for any
corporate group.

\begin{blockquote}
In Italy, where the majority of students were mature men pursuing advanced study in law and
medicine, their guilds came to exercise great power. Students hired and paid teachers, determined
the courses to be given, and fined any lecturer who skipped a chapter in expounding his subject.
\end{blockquote}

At French and English universities, where students were younger, masters’ guilds had the upper hand.
They forbade students to swear or gamble, fined them for breaking curfews, and prescribed table
manners. “Cleanse not the teeth,” one rule went, “with the steel that is sharpened for those that
eat with thee.”

Upon hearing “university” we tend to envision ivied halls and grassy quadrangles. Medieval
universities, however, had not even the semblance of permanent quarters. At first lectures were
given in wayside sheds at Oxford and Cambridge, in the cathedral cloisters in Paris, and in the
squares in Italy. In time teachers rented rooms. Students sat on the floor, which was usually
covered with straw to protect against the dampness. Unencumbered with athletic stadiums, libraries,
or other equipment, universities could pick up and move elsewhere at any time if they found
themselves at odds with local citizens.

In addition to lectures, the method of teaching was the disputation. Two or more masters, and
occasionally the students, debated text readings, employing Abelard’s question-and-answer approach.
Scholasticism developed in this context and came to stand for painstaking arrival at logical
conclusions through questioning, examining, and arranging details into a system of logic. The
Scholastic disputation stirred heated clashes and bitter feelings. Wars of logic ran for years
between master and master, with adherents of each cheering their hero on with tumultuous stomping
and whistling. Something important was happening in this raucous atmosphere. Students were learning
to think. Unquestioning acceptance of traditional authorities was no longer assured; even the
conclusions of Christian doctrine were being investigated.

\section{SUPPORT FOR THE PAPAL MONARCHY}

All of this feverish activity proceeded, however, under the watchful eye of the papacy. Some debates
continued for generations, but the popes worked to make sure that the net result was a new
philosophical framework that supported the papal monarchy. On the one hand was a new formulation of
canon law, and on the other hand a systematic statement of Christian theology.

The University of Bologna in Italy emerged in the later years of the twelfth century as the center
of the study of Roman civil law and church or, as it is called, canon law.

Canon law served the church just as civil law supported a secular government. It defined the rights,
duties, and powers of all people and priests within the church. It was the law administered in all
ecclesiastical courts, from those of the bishop up to that of the pope.

Sometime around 1140 Gratian, a monk of the monastery of St. Felix at Bologna, published a Harmony
of Discordant Canons, which tried to coordinate all previous collections of church law. Since he
arranged his quotations of authorities subject by subject, his Harmony soon emerged as the sole
manual for teachers and for judges in the church.

It served as the base for later additions to the Body of Canon Law. Thus by the fourteenth century
the Roman Church had at its disposal—until 1918, when it was revised—an authoritative body of laws
to control and direct the lives of men.

Canon law interfered in wars and recognized the care of widows. It decreed fasting and feasting. It
ordained confession and communion at least once a year. It specified which actions excommunicated a
Christian from the church.

But it also directed man’s most intimate relationships. By virtue of its concern with baptism it
established standards for all births—and all that led to births. The first inviting smile between
man and woman brought the couple under its watchful shadow. It directed penance for fornication and
adultery; it laid down the conditions in which a marriage could exist.

Canon law, in short, reached not only to every priest, but to every layman, plowman, or prince. And
it professed to declare not only the necessary path to salvation but the very nature of the most
intimate organs of men and women.

This craving to control, the zeal to reach every aspect of life, may reflect a Christian legalism.
It tried to reduce human freedom to a miserable minimum and forgot that God designed freedom as the
forerunner for the coming of true faith.

Canon law, however, gave to the papacy a rational legal basis, something the medieval state did not
yet possess. As a result the papacy rose to preeminent power in the public life of Europe and
achieved an international prestige that far outweighed that of any feudal kingdom.

On the day of Innocent III’s consecration to the papal office, he took as the text of his sermon
God’s words in Jeremiah 1:10: “I have this day set thee over the nations and over the kingdoms, to
root out, and to pull down, to destroy and to throw down, to build, and to plant.” For three decades
that is what he did. He was the “Vicar of Christ” on earth, thanks in large part to the universal
authority conveyed in canon law.

The second way the universities served the universal papacy was to provide an unshakable, rational
theological construction of Christian society. By the thirteenth century, the universities faced a
controversial issue. Several writings of Aristotle had been re-introduced to the West in the new
universities, lost works that had survived the decline of Rome among the eastern Christians, and
especially among the Muslims. The prospect of employing Aristotelian philosophy, however, raised
challenges. For example, the new Aristotelian approach called for a new matrix or framework, a
revised worldview, whereas most theologians had previously assumed a Platonic frame of reference.

A basic dualism followed: theologians distinguished between the basic material stuff of the world we
live in and the mental or intellectual realm (Plato’s world of Forms). The conceptual realm held
basic ideas, everything from the idea of a circle (circularity) to the idea of justice in and of
itself. The ideas are perfect, non-material, and eternal. Our concrete world is composed of
imperfect, material, and corruptible (not eternal) copies of these perfect ideas. This world
contains many circles, but every physical, tangible circle is imperfect (out of round) and temporal
(subject to time and its effects; not eternal).

Theologians such as Origen and Augustine had told the Christian story with this philosophical
backdrop in place. One might picture God and his ideas in an eternal realm and the corrupt and
compromised physical world below. While Aristotle did not overturn everything about Platonism, he
thought that the ideas or forms were present in the concrete physical realities. Discovering the
forms then called for one to investigate the concrete world or tangible creation. The overall thrust
for Aristotle was not to separate two unfitting dimensions (the conceptual and material), but was to
discover and admire how the best thinking integrated or fit with the material world by God’s special
design. This shift in viewing the world and reason was challenging in its scale, its novelty, and
its uncertain future.

Many scholars only knew Aristotle as he was presented by a famous Muslim interpreter named Averroes.
He stressed philosophy or reason’s independence from revelation or faith. His work left the
impression that one might follow reason to one conclusion, while revelation might point the learner
in a completely different direction. At times, Averroes would follow Aristotle and reason to
conclusions contradicting Christian truth, such as matter being eternal (not created). Things
changed with the arrival of Thomas Aquinas. His life work was to insist that good thinking never
leads the learner away from God’s truth. The truth that God reveals surpasses the truth we can
achieve by reasoning, but they never contradict. As the popular saying goes, “All truth is God’s
truth.”

\section{THE PEAK OF RATIONAL REACH}

To address controversies concerning doctrine, Thomas Aquinas (1224–1274) was dispatched to Paris
from Italy. Thomas Aquinas (for his father’s home, Aquino) was a Dominican monk of noble birth,
brilliant mind, tireless industry, and gentle disposition: he, too, had turned to Aristotle but he
had distinguished himself for his fidelity to the church. Unlike Averroes, Thomas taught that reason
fit with revelation. The result was his Summa Theologica (a summation of theological knowledge).

The Summa has in view the whole universe. Thomas says at the beginning, “In sacred doctrine
(theology) all things are treated from the standpoint of God, and the content of theology is in part
God Himself and in part other beings in as far as they are ordained unto God.”

Aquinas made a clear distinction between philosophy and theology, reason and revelation, but there
is no contradiction between the two. Both are fountains of knowledge; both come from the same God.

\figcaption{Thomas Aquinas (1224–1274) formulated Christian doctrine using the rational method. His
masterpiece, Summa Theologica, is still widely read by theologians and philosophers.}

The two differ in their methods of searching after truth. Reason is based upon the visible creation
and can reach ideas that deal with “the vestibule of faith.” Revelation looks to God as he is in
himself and so is superior to reason both in its certainty and in its subject matter.

Reason, for example, can prove God’s existence. Accepting Aristotle’s principle (every effect has a
cause, every cause a prior cause, and so on back to the First Cause) Thomas declared that creation
traces back to a divine First Cause, the Creator.

However, the full knowledge of God—the Trinity, for example—comes only through revelation. From this
knowledge we discover man’s origin and destiny.

Man is a sinner and in need of special grace from God. Jesus Christ, by his sacrifice, has secured
the reconciliation of man and God. All who receive the benefits of Christ’s work are justified, but
the key, as in traditional Catholic teaching, lies in the way the benefits of Christ’s work are
applied. Christ won grace; the church imparts it. Aquinas taught that Christians need the constant
infusion of “cooperating grace,” whereby the Christian virtues—above all, love—are stimulated in the
soul. Assisted by this cooperating grace, a Christian can do works that please God and gain special
merit in God’s sight.

This saving grace, said Aquinas, comes to men exclusively through the channel of divinely appointed
sacraments placed in the keeping of the church, the visible, organized Roman body, led by the pope.
So convinced was Aquinas of the divine sanction of the papacy that he insisted that submission to
the pope was necessary for salvation.

Aquinas followed Peter Lombard who had written the standard textbook for theology, in holding to
seven sacraments: baptism, confirmation, the Lord’s Supper, penance, extreme unction, marriage, and
ordination.

The sacrament of all sacraments, however, is the Lord’s Supper, which is more than the communion of
the early church. As the Roman Church had for centuries, Aquinas held that it is a true sacrifice,
continuing that of Christ on the cross, and predisposing God to be gracious to those for whom it was
offered. In the Supper the essence or genuine being of the bread and wine are changed miraculously
into the actual body and blood of Christ while the exterior remains unchanged. Thomas gives the
classic presentation of this doctrine known as transubstantiation. Since sin remains a problem for
the baptized Christian, God provided penance, the sacrament of spiritual healing. According to
Thomas it consists of three elements: contrition or sorrow for the sin; confession to the priest as
the spiritual physician who can apply the appropriate remedy and pronounce absolution; and
satisfaction, by which the Christian makes good the evil effects of the sin, usually called an “act
of penance.”

With some cautions Thomas also accepted the practice of indulgences that had gained prominence
during the crusades. Aquinas taught that thanks to the work of Christ and the meritorious deeds of
the saints, the church has access to a “treasury of merit,” a great spiritual reservoir. Priests may
draw from this to aid Christians who have insufficient merit of their own.

Death is the great divide. The wicked, said Thomas, pass into hell. The faithful who have wisely
used the means of grace pass immediately to heaven. But the mass of mankind, who while Christian in
desire and participants in the sacraments, have followed Christ inadequately, must suffer further
purification in purgatory before attaining the joys of heaven. Thankfully these souls are not beyond
the help of the church on earth. Prayers to the saints in heaven can prevail to relieve the pains of
souls in purgatory.

Thus we come to the peak of ecclesiastical aspiration—and arrogance. Earth alone is not enough! The
pope and his priests not only mediate the grace of God to sinners on earth by the miracle of the
Blessed Sacrifice and by their prayers for the dead; they reach beyond the grave to minister to
suffering souls.

There is nothing new in this. It had been said many times before. But Thomas set the traditional
teachings of the church in a grand, almost cosmic, framework. Like the Gothic cathedral, his system
aimed at a perfect harmony between the aspirations of man and the light of God’s truth. It provided
the papal monarchy an impressive supernatural view of this world and the next.

But like the crusades, Scholastic theology may have reached too far, claimed too much for itself and
the church. A clue to the heights of arrogance lies in Innocent III’s claim that the pope is the
judge of the world, “set in the midst between God and man, below God and above man.” The boast did
not go unchallenged. A growing host of unlettered laymen remembered the apostolic testimony, “There
is one God and one Mediator.”

\readinghead{Suggestions for Further Reading}
\begin{readinglist}
\entry *Barron, Robert. Thomas Aquinas: Spiritual Master. 2nd ed. New York: Crossroad, 2008.
\entry *Bauerschmidt, Frederick. Holy Teaching: Introducing the Summa Theologiae of St. Thomas.
  Grand Rapids: Brazos, 2005.
\entry Daniel-Rops, H. Cathedral and Crusade. London: J. M. Dent, 1957.
\entry Dauphinais, Michael, and Matthew Levering. Knowing the Love of Christ: An Introduction to the
  Theology of St. Thomas Aquinas. Notre Dame: University of Notre Dame Press, 2002.
\entry Fremantle, Anne. Age of Faith. New York: Time-Life Books, 1968.
\entry Haskins, Charles Homer. The Rise of the Universities. Ithaca, NY: Cornell University Press,
  1957.
\entry *Healy, Nicholas. Thomas Aquinas: Theologian of the Christian Life. Burlington, VT: Ashgate,
  2003.
\entry McGiffert, Arthur C. A History of Christian Thought: From Tertullian to Erasmus. New York:
  Scribner’s, 1954.
\entry Nichols, Aidan. Discovering Aquinas: An Introduction to His Life, Work, and Influence. Grand
  Rapids: Eerdmans, 2003.
\end{readinglist}
